\label{chap:83-18}

The inverse limit \(\Omega_\infty\) preserves the complete history. To
obtain a real coordinate, a compatible family of intervals must be
added and three different facts proved: each history determines a
unique value; every value of the declared compact set is attained; and the finite
measures descend consistently. None of these facts follows
from the mere growth in the number of states. The construction below
proves all three facts jointly for the HMT system.

The Archimedean publication likewise preserves an approximation of the
identity by cylindrical conditional expectations. The concentration of unit
mass on compatible prefixes provides the exact passage from the
nonadic cylinders to the Dirac measure, without identifying density,
support measure, and total mass.

\begin{HMTScientificOwnerSubordinated}
\input{ampliaciones_sucesoras_20260824/parte_i_ii/owners/03_cilindros_dirac}
\end{HMTScientificOwnerSubordinated}

\section{System of intervals associated with finite states}

Let \((S_n,r_{n,m})\) be the inverse system of the
\cref{chap:sistema-inverso-continuidad-genealogica}. For each
\(a\in S_n\), assign a nonempty compact interval
\begin{equation}
 I_n(a)=[\ell_n(a),u_n(a)]\subset\mathbb R.
 \label{p12-83-c18-s1:eq:intervalo-asociado-estado}
\end{equation}
The assignment is \emph{compatible with refinement} when
\begin{equation}
 r_{n+1,n}(b)=a
 \quad\Longrightarrow\quad
 I_{n+1}(b)\subseteq I_n(a).
 \label{p12-83-c18-s1:eq:anidamiento-intervalos-estados}
\end{equation}
It is \emph{uniformly contractive} when there exists a sequence
\(\delta_n\downarrow0\) such that
\begin{equation}
 \operatorname{diam}I_n(a)\leq\delta_n
 \qquad(a\in S_n).
 \label{p12-83-c18-s1:eq:contraccion-uniforme-intervalos}
\end{equation}

For a history \(x=(x_n)_n\), the inclusions
\eqref{p12-83-c18-s1:eq:anidamiento-intervalos-estados} produce a chain
\begin{equation}
 I_0(x_0)\supseteq I_1(x_1)\supseteq I_2(x_2)\supseteq\cdots.
 \label{p12-83-c18-s1:eq:cadena-intervalos-historia}
\end{equation}

\begin{theorem}[Evaluation of a history]
\label{p12-83-c18-s1:thm:evaluacion-unica-historia-u024}
If the assignment is compatible and uniformly contractive, then for every \(x\in\Omega_\infty\) there exists a unique real number
\begin{equation}
 \operatorname{ev}_{\mathbb R}(x)
 \in\bigcap_{n\geq0}I_n(x_n).
 \label{p12-83-c18-s1:eq:evaluacion-real-historia-u024}
\end{equation}
The map
\begin{equation}
 \operatorname{ev}_{\mathbb R}:\Omega_\infty\longrightarrow\mathbb R
 \label{p12-83-c18-s1:eq:aplicacion-evaluacion-real-u024}
\end{equation}
is continuous.
\end{theorem}

\begin{proof}
The intervals of \eqref{p12-83-c18-s1:eq:cadena-intervalos-historia} are compact,
nonempty and nested. The nested-interval theorem guarantees that
their intersection is nonempty. If \(s,t\) belong to it, then
\[
 |s-t|\leq\operatorname{diam}I_n(x_n)\leq\delta_n
\]
for every \(n\). Setting \(n\to\infty\) yields \(s=t\).

For continuity, let \(\varepsilon>0\) and choose \(n\) with
\(\delta_n<\varepsilon\). If \(y\in[x_n]_n\), then
\(I_n(y_n)=I_n(x_n)\). The two evaluations belong to that interval and,
therefore,
\[
 |\operatorname{ev}_{\mathbb R}(x)
  -\operatorname{ev}_{\mathbb R}(y)|<\varepsilon.
\]
The cylinder \([x_n]_n\) is an open neighborhood of \(x\).
\end{proof}

The construction does not identify \(x\) with
\(\operatorname{ev}_{\mathbb R}(x)\). A single coordinate may have
several histories, and a history contains more data than any point on
the line.

\section{Coherent coverage of the visible compactum}

Let \(K\subset\mathbb R\) be a nonempty compact set. The HMT construction of the
cylinders associates with nesting and contraction the coverage
identities
\begin{align}
 K&=\bigcup_{a\in S_0}I_0(a),
 \label{p12-83-c18-s1:eq:cobertura-inicial-compacto-u024}\\
 I_n(a)&=
 \bigcup_{\substack{b\in S_{n+1}\\r_{n+1,n}(b)=a}}
 I_{n+1}(b)
 \qquad(a\in S_n).
 \label{p12-83-c18-s1:eq:cobertura-extensiones-compatible-u024}
\end{align}
The second equality is the geometric publication of the prolongation law:
each point of \(I_n(a)\) remains in a compatible child cylinder. Together with
the surjectivity of \(r_{n+1,n}\), already established by prefix
survival, it simultaneously produces states above \(a\) and the
complete coverage of \(I_n(a)\).

\begin{theorem}[Surjectivity of the evaluation]
\label{p12-83-c18-s1:thm:sobreyectividad-evaluacion-compacto}
The constructed identities
\eqref{p12-83-c18-s1:eq:cobertura-inicial-compacto-u024} and
\eqref{p12-83-c18-s1:eq:cobertura-extensiones-compatible-u024} imply
\begin{equation}
 \operatorname{ev}_{\mathbb R}(\Omega_\infty)=K.
 \label{p12-83-c18-s1:eq:imagen-evaluacion-compacto}
\end{equation}
\end{theorem}

\begin{proof}
The left-to-right inclusion follows from
\(I_n(x_n)\subseteq I_0(x_0)\subseteq K\).

Let us fix \(t\in K\). Let us define
\begin{equation}
 T_n(t):=\{a\in S_n:t\in I_n(a)\}.
 \label{p12-83-c18-s1:eq:estados-que-cubren-punto}
\end{equation}
The initial cover makes \(T_0(t)\) nonempty. If \(a\in T_n(t)\), the identity \eqref{p12-83-c18-s1:eq:cobertura-extensiones-compatible-u024} provides at least one \(b\in T_{n+1}(t)\) with \(r_{n+1,n}(b)=a\). The graph of these extensions is finitely branching and has nonempty levels at every depth. There therefore exists a compatible chain \(x=(x_n)_n\) with \(t\in I_n(x_n)\) for every \(n\). Uniqueness of the intersection forces \(\operatorname{ev}_{\mathbb R}(x)=t\).
\end{proof}

\section{Archimedean genealogical quotient and homeomorphism with \texorpdfstring{\(\mathbb R\)}{R}}

We define the relation
\begin{equation}
 x\sim_{\mathbb R}y
 \quad\Longleftrightarrow\quad
 \operatorname{ev}_{\mathbb R}(x)
 =\operatorname{ev}_{\mathbb R}(y).
 \label{p12-83-c18-s1:eq:relacion-equivalencia-arquimediana}
\end{equation}
Let
\begin{equation}
 q_{\mathbb R}:\Omega_\infty
 \longrightarrow\Omega_\infty/\!\sim_{\mathbb R}
 \label{p12-83-c18-s1:eq:proyeccion-cociente-arquimediano}
\end{equation}
the canonical projection.

\begin{theorem}[Reconstruction of the compact as a quotient]
\label{p12-83-c18-s1:thm:homeomorfismo-cociente-compacto-u024}
For the preceding HMT system there is a unique homeomorphism
\begin{equation}
 \overline{\operatorname{ev}}_{\mathbb R}:
 \Omega_\infty/\!\sim_{\mathbb R}
 \xrightarrow{\ \cong\ }K
 \label{p12-83-c18-s1:eq:homeomorfismo-cociente-compacto-u024}
\end{equation}
such that
\begin{equation}
 \operatorname{ev}_{\mathbb R}
 =\overline{\operatorname{ev}}_{\mathbb R}\circ q_{\mathbb R}.
 \label{p12-83-c18-s1:eq:factorizacion-evaluacion-cociente}
\end{equation}
The fiber over \(t\in K\) is
\begin{equation}
 \mathfrak G_t
 :=\operatorname{ev}_{\mathbb R}^{-1}(t),
 \label{p12-83-c18-s1:eq:fibra-genealogias-mismo-valor}
\end{equation}
the complete set of genealogies that publish \(t\).
\end{theorem}

\begin{proof}
The map \(\operatorname{ev}_{\mathbb R}\) is continuous, surjective, and
constant exactly on the classes of \(\sim_{\mathbb R}\). By the
universal property of the quotient, it induces a continuous bijective map
\(\overline{\operatorname{ev}}_{\mathbb R}\). The domain is compact because it is a continuous
image of the compact space \(\Omega_\infty\), and \(K\) is Hausdorff.
A continuous bijection from a compact space to a Hausdorff space is a homeomorphism.
\end{proof}

The theorem proves two things and keeps them separate:
\begin{equation}
 \boxed{
 \begin{array}{c}
 \text{the history space is profinite and retains memory,}\\
 \text{its quotient by equality of evaluation is the visible compact space.}
 \end{array}}
 \label{p12-83-c18-s1:eq:dos-niveles-continuo-u024}
\end{equation}
It is not asserted that \(\Omega_\infty\) is homeomorphic to \(K\).

\section{Genealogical Atlas of the Real Line}

For \(m\geq1\), let \(K_m=[-m,m]\). The same construction, applied to each
\(K_m\), produces a system \(\Omega_\infty^{(m)}\) that reconstructs
\(K_m\), together with closed inclusions
\begin{equation}
 j_m:\Omega_\infty^{(m)}\hookrightarrow
 \Omega_\infty^{(m+1)}
 \label{p12-83-c18-s1:eq:inclusion-atlas-genealogico}
\end{equation}
which satisfy
\begin{equation}
 \operatorname{ev}_{m+1}\circ j_m
 =\iota_m\circ\operatorname{ev}_m,
 \label{p12-83-c18-s1:eq:compatibilidad-atlas-real}
\end{equation}
where \(\iota_m:K_m\hookrightarrow K_{m+1}\) is the inclusion ordinary.

\begin{proposition}[Exhaustion of the line]
\label{p12-83-c18-s1:prop:agotamiento-recta-real-u024}
The Archimedean quotient of the strict colimit
\begin{equation}
 \Omega_\infty^{\mathbb R}
 :=\varinjlim_m\Omega_\infty^{(m)}
 \label{p12-83-c18-s1:eq:colimite-historias-recta}
\end{equation}
is
\begin{equation}
 \varinjlim_m[-m,m]
 =\bigcup_{m\geq1}[-m,m]
 =\mathbb R
 \label{p12-83-c18-s1:eq:colimite-compactos-recta}
\end{equation}
with its usual topology.
\end{proposition}

\begin{proof}
Each piece has quotient \(K_m\) by the
\cref{p12-83-c18-s1:thm:homeomorfismo-cociente-compacto-u024};
\eqref{p12-83-c18-s1:eq:compatibilidad-atlas-real} intertwines the quotients. The final topology
of the exhaustion by compact sets coincides with the usual topology of
\(\mathbb R\): a subset \(U\subseteq\mathbb R\) is open if and only
if \(U\cap[-m,m]\) is open in \([-m,m]\) for every \(m\).
\end{proof}

Passage from a bounded interval to the line does not erase the fibers
\(\mathfrak G_t\). The inclusions \(j_m\) preserve the history and only
enlarge the available visible chart.

\section{Balanced ternary and block-decimal charts}

The balanced ternary representation of an integer has the form
\begin{equation}
 n=\sum_{j=0}^{J}\epsilon_j3^j,
 \qquad
 \epsilon_j\in\{-1,0,+1\},
 \label{p12-83-c18-s1:eq:carta-ternaria-balanceada-u024}
\end{equation}
with a fixed carry rule. This chart preserves the local sign and the
quotient of the reduction; it must not be confused with the nonnegative word in
\(\mathbb F_3\).

The block decimal chart utilizes
\begin{equation}
 D_N=1000D_{N-1}+u_N,
 \qquad
 u_N\in\{0,\ldots,999\},
 \qquad D_0=0.
 \label{p12-83-c18-s1:eq:concatenacion-base-mil-u024}
\end{equation}
The block \(u_N\), its class modulo nine, and the carry are distinct fields.
Since \(1000\equiv1\pmod9\),
\begin{equation}
 D_N\equiv\sum_{j=1}^{N}u_j\pmod9.
 \label{p12-83-c18-s1:eq:base-mil-compatible-mod9-u024}
\end{equation}
This congruence does not recover the order of the blocks; the sequence
\((u_1,\ldots,u_N)\) remains in the record.

For \(M\geq1\), the cells decimales half-open are
\begin{equation}
 \mathcal D_{M,q}
 :=[q10^{-M},(q+1)10^{-M}),
 \qquad q\in\mathbb Z.
 \label{p12-83-c18-s1:eq:celda-decimal-semiabierta-u024}
\end{equation}
A half-open interval \([\ell,h)\) is contained in a unique cell if and only if
\begin{equation}
 \left\lfloor10^M\ell\right\rfloor
 =\left\lceil10^Mh\right\rceil-1.
 \label{p12-83-c18-s1:eq:criterio-exacto-celda-decimal-u024}
\end{equation}
The formula with a ceiling at the right endpoint also covers the case in which
\(h\) coincides exactly with a decimal boundary.

\begin{definition}[Truncation and Rounding]
For \(x\geq0\), define
\begin{align}
 \operatorname{Tr}_M(x)
 &:=10^{-M}\left\lfloor10^Mx\right\rfloor,
 \label{p12-83-c18-s1:eq:truncamiento-decimal-u024}\\
 \operatorname{Rd}_M(x)
 &:=10^{-M}\left\lfloor10^Mx+\frac12\right\rfloor.
 \label{p12-83-c18-s1:eq:redondeo-decimal-u024}
\end{align}
\end{definition}
These operations coincide except when the discarded fraction crosses the
half-unit threshold at the final position. The distinction is semantic and
arithmetic; one cannot be substituted for the other to force prefix
compatibility.

\section{Descent of operations along evaluation fibers}

Let \(D_\star\subseteq\Omega_\infty\times\Omega_\infty\) the domain of to
operation parcial continuous
\begin{equation}
 \star:D_\star\longrightarrow\Omega_\infty.
 \label{p12-83-c18-s1:eq:operacion-historias-u024}
\end{equation}
Assume that \(D_\star\) is saturated by the fibers of
\(\operatorname{ev}_{\mathbb R}\times\operatorname{ev}_{\mathbb R}\).

\begin{theorem}[Exact Descent Criterion]
\label{p12-83-c18-s1:thm:criterio-descenso-operaciones-u024}
There exists a unique visible operation
\begin{equation}
 \overline\star:
 (\operatorname{ev}_{\mathbb R}\times\operatorname{ev}_{\mathbb R})
 (D_\star)\longrightarrow K
 \label{p12-83-c18-s1:eq:operacion-descendida-visible}
\end{equation}
which satisfies
\begin{equation}
 \operatorname{ev}_{\mathbb R}(x\star y)
 =\operatorname{ev}_{\mathbb R}(x)
  \ \overline\star\ 
  \operatorname{ev}_{\mathbb R}(y)
 \label{p12-83-c18-s1:eq:entrelazamiento-operacion-descendida}
\end{equation}
if and only if
\begin{equation}
 \begin{split}
 &\operatorname{ev}_{\mathbb R}(x)=
  \operatorname{ev}_{\mathbb R}(x'),\\
 &\operatorname{ev}_{\mathbb R}(y)=
  \operatorname{ev}_{\mathbb R}(y')
 \end{split}
 \quad\Longrightarrow\quad
 \operatorname{ev}_{\mathbb R}(x\star y)=
 \operatorname{ev}_{\mathbb R}(x'\star y')
 \label{p12-83-c18-s1:eq:constancia-operacion-sobre-fibras}
\end{equation}
for all pairs in the domain.
\end{theorem}

\begin{proof}
If \(\overline\star\) exists, the right-hand side of
\eqref{p12-83-c18-s1:eq:entrelazamiento-operacion-descendida} depends only on the two
evaluations, and
\eqref{p12-83-c18-s1:eq:constancia-operacion-sobre-fibras} follows. Conversely, this
constancy makes it possible to define
\[
 s\ \overline\star\ t
 :=\operatorname{ev}_{\mathbb R}(x\star y)
\]
for any representatives \(x,y\) with evaluations \(s,t\). The
condition proves that the definition is independent of the representatives.
Uniqueness is immediate.
\end{proof}

The additive and multiplicative sheets of APP supply operations on a
common support, but their carries and their memories are not identified. The
preceding theorem allows sum and product to descend to the visible chart
without requiring their genealogical realizations to coincide.

\section{Correspondence between discrete cardinality, cylindrical measure, and Archimedean publication}

At each depth three operations appear that must be kept separated:
\begin{align}
 \operatorname{Count}_n(A)&:=|A|,
 \label{p12-83-c18-s1:eq:operacion-conteo-u024}\\
 \operatorname{Measure}_n(A)&:=\mu_n(A),
 \label{p12-83-c18-s1:eq:operacion-medida-u024}\\
 \operatorname{Publish}_n(a)&:=L_n(a).
 \label{p12-83-c18-s1:eq:operacion-publicacion-u024}
\end{align}
The count depends on the number of states; the measure depends on compatible
weights; the publication depends on the reader. Two fibers of equal
cardinality can have distinct weights, and two states of equal publication
can belong to distinct cylinders. Equality of one of these three
outputs does not imply equality of the other two.

\section{Obstructions to Archimedean evaluation: nesting, continuity, quotient, projective measures, and descent of operations}

\begin{criterion}[Refutation of the evaluation and the measure]
The construction is refuted if any of the following facts
holds:
\begin{enumerate}
 \item a compatible extension receives an interval that is not contained
       in the interval of its restriction;
 \item the diameters do not tend to zero, but uniqueness of the
       intersection is asserted;
 \item the evaluation is not continuous in the cylinder topology;
 \item it is asserted that every point of \(K\) appears without proving the coherent
       covering \eqref{p12-83-c18-s1:eq:cobertura-extensiones-compatible-u024};
 \item the map induced by the quotient is not bijective;
 \item the profinite space is identified with its visible quotient;
 \item the inclusions of the atlas do not intertwine their evaluations;
 \item \(\lfloor10^Mh\rfloor\) is used instead of
       \(\lceil10^Mh\rceil-1\) for a semi-open interval, and the
       boundary case is lost;
 \item truncation and rounding are conflated;
 \item a descended operation depends on the representative chosen in a
       fibre;
 \item to family \((\mu_n)_n\) viola
       \eqref{eq:consistencia-proyectiva-medidas-u024};
 \item the measure of a cylinder depends on the level to which it is refined;
 \item a conditioned measure ceases to sum to one or loses consistency;
 \item the visible measure \(\nu\) is identified with the internal measures
       \(\mu_t\);
 \item an arrow between \(\mathbb R\) and
       \(\mathcal E_\infty\) is inferred solely because both publications share an
       antecedent;
 \item the count of a fibre is used as though it determined its measure or its
       published value.
\end{enumerate}
\end{criterion}


\section{From the space of histories to the real line}

Let
\[
 \Psi:\mathbb Z_3^6\xrightarrow{\ \sim\ }
       (\mathbb F_3^6)^{\mathbb N}
\]
the Hensel conjugacy between the initial memory and the complete tape of
six-trit blocks. Concatenating the six coordinates of each block yields
a ternary tape
\(
 (a_n)_{n\geq 1}\in\{0,1,2\}^{\mathbb N}
\), on which the evaluation acts
\[
 \operatorname{ev}_3((a_n))
 :=\sum_{n\geq1}\frac{a_n}{3^n}\in[0,1].
\]
The two expansions of a ternary endpoint represent the same value and
remain distinct before evaluation is applied. This does not constitute
an ambiguity of the genealogical object: it is precisely a fiber of the reader.

We define
\[
 X_{\mathrm{raw}}:=\mathbb Z\times\mathbb Z_3^6,
 \qquad
 P(m,x):=m+\operatorname{ev}_3(\operatorname{flat}(\Psi x)).
\]
The integer coordinate provides the countable atlas of the line, and the
3-adic coordinate preserves the complete history of the fractional part.

\begin{theorem}[Exact reconstruction of the Archimedean quotient]
\label{p12-83-c18-s2:hmtch:thm-cociente-real}
The map \(P:X_{\mathrm{raw}}\to\mathbb R\) is sobreyectiva. If
\[
 x\sim_P y
 \quad\Longleftrightarrow\quad
 P(x)=P(y),
\]
then
\[
 \overline P:X_{\mathrm{raw}}/\!\sim_P\;\longrightarrow\mathbb R,
 \qquad [x]\longmapsto P(x),
\]
is a bijection. Consequently,
\[
 \left|X_{\mathrm{raw}}/\!\sim_P\right|
 =|\mathbb R|
 =2^{\aleph_0}.
\]
\end{theorem}

\begin{proof}
Every ternary sequence is grouped uniquely into consecutive blocks of
six digits. The conjugation \(\Psi\) assigns a unique initial
3-adic memory to that tape. Given \(r\in\mathbb R\), take
\(m=\lfloor r\rfloor\) and a ternary expansion of
\(r-m\in[0,1)\); the preceding reconstruction produces
\(x\in\mathbb Z_3^6\) with \(P(m,x)=r\). Thus, \(P\) is surjective.

The map \(\overline P\) is well defined by the definition of
\(\sim_P\), is surjective because \(P\) is, and is injective because two
classes with the same image belong to the same fiber. Finally,
\[
 |\mathbb Z_3^6|
 =\left|(\mathbb F_3^6)^{\mathbb N}\right|
 =729^{\aleph_0}
 =2^{\aleph_0},
\]
and multiplying by the countable coordinate \(\mathbb Z\) does not alter that
cardinal. The calculation agrees with that of the quotient, already identified with
\(\mathbb R\).
\end{proof}

\begin{remark}[Surviving closure]
\label{p12-83-c18-s2:hmtch:rem-version-superviviente}
The raw completion and the surviving subspace
\(X_{\mathrm{surv}}\) are two stages of the same construction. The survival
law produces coherent coverings at each level, uniform contraction
of diameters and a prolongation for every visible cell. The infinite-branch
lemma then produces a section over each point, and the
evaluation induces, compact set by compact set, a homeomorphism of the quotient with the
published interval. The exhaustion by \([-m,m]\) reconstructs the entire line,
in agreement with
\cref{p12-83-c18-s1:prop:agotamiento-recta-real-u024,p12-83-c18-s2:hmtch:thm-cociente-real}.
\end{remark}

\section{Archimedean mean and double projection of the same state}

Let \(u=(m,\mathbf u)\in
\mathscr O_{\mathbb R,\leftrightarrow}^{\rm HMT}\). A prefix
\(\mathbf u_n\) determines in the fractional chart an oriented
semi-open cylinder
\[
 I_n(\mathbf u_n)=[a_n,b_n),
\]
with the complementary convention at the endpoints of double expansion. To
apply compactness, its closure
\(K_n=[a_n,b_n]\) is used. The compatibilities give
\[
 K_{n+1}\subseteq K_n,
 \qquad
 \delta_n:=b_n-a_n\longrightarrow0.
\]
The global Archimedean mean of the cylinder is
\[
 \mu_n(u)=m+\frac{a_n+b_n}{2}.
\]
If \(q\geq n\), then
\[
 |\mu_q(u)-\mu_n(u)|\leq\delta_n,
\]
so that
\begin{equation}
 \mathfrak R(u):=\lim_{n\to\infty}\mu_n(u)
 \label{p12-83-c18-s2:hmtch:eq-media-arquimediana}
\end{equation}
exists. What tends to zero is the diameter of the cylinder; the real value
is the resultant of the nesting, not that which vanishes.

\begin{theorem}[Two-way quotient and real line]
\label{p12-83-c18-s2:hmtch:thm-cociente-two-way-real}
Let
\[
 u\sim_{\mathfrak R}v
 \quad\Longleftrightarrow\quad
 \mathfrak R(u)=\mathfrak R(v).
\]
Under the coherent covering of the HMT charts, the evaluation induces a homeomorphism.
\begin{equation}
 \overline{\mathfrak R}:
 \mathscr O_{\mathbb R,\leftrightarrow}^{\rm HMT}/\!
 \sim_{\mathfrak R}
 \xrightarrow{\ \sim\ }\mathbb R.
 \label{p12-83-c18-s2:hmtch:eq-homeomorfismo-two-way-real}
\end{equation}
In particular, this is a continuous HMT morphism: not the result of choosing
an isolated history from each fiber.
\end{theorem}

\begin{proof}
The cylinder bound proves the continuity of \(\mathfrak R\). The covering produces a history over every real and hence surjectivity. Upon passage to the quotient by the fibers, the induced map is bijective. On every compact chart, a continuous bijection onto a Hausdorff interval is a homeomorphism. The global topology of the quotient is the final topology of the gluing of these charts; the evaluation fiber identifies endpoint \(m+1\) of one chart with the translated endpoint \(0\) of the following chart. The resulting cover is locally finite, so the chart homeomorphisms glue into a global homeomorphism onto \(\mathbb R\).
\end{proof}
